% ===========================================================================
%  The Singular Value Decomposition -- a Loom fill-in notebook.
%  Original exposition of standard (public-domain) linear algebra.
%  Compile:  latexmk -xelatex main      (or: xelatex main, twice)
% ===========================================================================
\documentclass{loom}

% per-document math shorthands
\newcommand{\R}{\mathbb{R}}
\newcommand{\C}{\mathbb{C}}
\newcommand{\T}{^{\!\top}}
\newcommand{\norm}[1]{\lVert #1\rVert}
\newcommand{\Sig}{\Sigma}
\DeclareMathOperator{\rank}{rank}
\DeclareMathOperator{\col}{col}
\DeclareMathOperator{\row}{row}
\DeclareMathOperator{\Null}{null}
\DeclareMathOperator{\tr}{tr}
\DeclareMathOperator{\diag}{diag}

\runningthread{singular value decomposition}

\begin{document}

\loomcover
  {The Singular Value Decomposition}
  {One decomposition, every question · fill-in notes}
  {woven on \textsc{Loom}}
  {\today}

% ===========================================================================
\section*{The one idea, and how to read these notes}
\addcontentsline{toc}{section}{The one idea · how to read}
\warmth{0}

\begin{strand}
Every real matrix $A$ factors as
\[
  A \;=\; U\,\Sig\,V\T,\qquad U,V \text{ orthogonal},\quad
  \Sig=\diag(\sigma_1\ge\sigma_2\ge\cdots\ge0).
\]
The numbers $\sigma_i$ are the \keyword{singular values}. The single idea of these
notes: \keyword{ask a matrix almost any question, and the answer is read straight off
its singular values.} Rank, size, invertibility, the best low-rank picture, the four
subspaces, the condition number: all of them live in $\Sig$.
\end{strand}

\block{The master table: read it off $\Sig$ (fill it in as you go)}
\noindent\begin{tabularx}{\linewidth}{@{}L L@{}}
\toprule
{\color{indigo}Question about $A$} & {\color{madder}Answer, read off $\Sig$}\\
\midrule
rank of $A$                         & number of $\fillin[2.5cm]$ singular values $=r$ \;(§\ref{sec:subspaces})\\
operator norm $\norm{A}_2$          & $\sigma\fillin[1cm]$ (the largest)\\
Frobenius norm $\norm{A}_F$         & $\sqrt{\sigma_1^2+\cdots+\sigma_r^2}$\\
is $A$ (square) invertible?         & yes iff $\sigma_n\,\fillin[1.2cm]\,0$\\
condition number $\kappa(A)$        & $\sigma_1/\sigma_r$ \;(§\ref{sec:faces})\\
best rank-$k$ approximation         & keep $\sigma_1,\dots,\sigma_k$, drop the rest \;(§\ref{sec:eckart})\\
$|\det A|$ (square)                 & $\sigma_1\sigma_2\cdots\sigma_n$\\
the four fundamental subspaces      & spans of columns of $U$ and $V$ \;(§\ref{sec:subspaces})\\
\bottomrule
\end{tabularx}

\block{How to read (passive \emph{and} active)}
\begin{itemize}[leftmargin=1.3em,itemsep=1pt]
  \item \keyword{Read} the strands and the stated theorems.
  \item \keyword{Fill} the blanks \fillin[1.2cm], the \TODO{missing step} proof gaps, and
        the \emph{Your turn} boxes. Your linear-algebra course is the answer key.
  \item Bump each section's \texttt{\textbackslash warmth\{0\}} to $0$--$5$ once you can rebuild it cold.
\end{itemize}

\bigskip
{\small\tableofcontents}

\clearpage
% ===========================================================================
% ===========================================================================
\section{The shape of a matrix}\label{sec:geometry}
\warmth{0}\quad\whisper{Every linear map is a rotation, a stretch along axes, and another rotation.}

\begin{strand}
Forget the entries for a moment and watch what $A\colon\R^n\to\R^m$ \emph{does}. Feed it the
unit sphere; out comes an \keyword{ellipsoid}. The singular value decomposition is just the
honest description of that ellipsoid: the directions of its axes, and how long they are.
Reading $A=U\Sig V\T$ right to left, $A$ is a \keyword{rotation} $V\T$, then an axis-aligned
\keyword{stretch} $\Sig$, then another rotation $U$. Three simple motions, in that order, are
\emph{all} any linear map ever does.
\end{strand}

\block{The decomposition}

\begin{definition}[singular value decomposition]\label{def:svd}
A \keyword{singular value decomposition} of $A\in\R^{m\times n}$ is a factorization
$A=U\Sig V\T$ where
\begin{itemize}[leftmargin=2em,itemsep=1pt]
  \item $U\in\R^{m\times m}$ and $V\in\R^{n\times n}$ are \keyword{orthogonal}:
        $U\T U=\fillin[1.5cm]$ and $V\T V=\fillin[1.5cm]$;
  \item $\Sig\in\R^{m\times n}$ is ``diagonal'' with entries
        $\sigma_1\ge\sigma_2\ge\cdots\ge\sigma_r>0=\sigma_{r+1}=\cdots$, the \keyword{singular values}.
\end{itemize}
The columns $u_i$ of $U$ are the \keyword{left} singular vectors; the columns $v_i$ of $V$ the
\keyword{right} singular vectors.
\end{definition}

\recall{Orthogonal matrices are exactly the rigid motions fixing the origin: they preserve every length and angle. $\det=\pm1$ (rotation or reflection).}

\block{The singular triple}
Columns of $A=U\Sig V\T$ say it most vividly. Each triple $(\sigma_i,u_i,v_i)$ satisfies
\[
  A\,v_i \;=\; \sigma_i\,u_i \qquad\text{and}\qquad A\T u_i \;=\; \sigma_i\,v_i .
\]
So $A$ takes the orthonormal directions $v_i$ to the orthogonal directions $\sigma_i u_i$:
the $v_i$ are the \emph{axes of the input circle}, the $\sigma_i u_i$ the \emph{axes of the
output ellipse}.

\begin{yourturn}
A maps the unit sphere to an ellipsoid whose semi-axis lengths are the singular values.
\begin{itemize}[leftmargin=1.3em,itemsep=1pt]
  \item For $A=\diag(3,2)$ the unit circle becomes an ellipse with semi-axes $\fillin[1cm]$ and
        $\fillin[1cm]$, along the $\fillin[2.5cm]$ directions.
  \item If some $\sigma_i=0$, the ellipsoid is \emph{flat} in that direction. Sketch what the
        unit circle maps to when $A=\diag(2,0)$. \recall{A flat ellipse = a line segment. That collapse is exactly rank loss; see §\ref{sec:subspaces}.}
\end{itemize}
\workspace[2]
\end{yourturn}

\begin{strand}
Why insist on the ordering $\sigma_1\ge\sigma_2\ge\cdots$? Because then the \emph{first} singular
triple is the single most important thing $A$ does: $\sigma_1$ is how much $A$ can stretch any
unit vector, achieved in direction $v_1$. Every later triple is the best that remains after the
earlier ones are used up. This ``greedy from the top'' structure is the whole reason truncating
the SVD gives the best simple picture (§\ref{sec:eckart}).
\end{strand}

\loose{Complex case: for $A\in\C^{m\times n}$, replace ``orthogonal/transpose'' by ``unitary/conjugate-transpose'' and every word here still holds. Singular values stay real and $\ge0$.}

% ===========================================================================
\section{Why it always exists}\label{sec:existence}
\warmth{0}\quad\whisper{Eigen-decomposition can fail; the SVD never does. The reason is $A\T A$.}

\begin{strand}
Eigenvalues are fussy: a matrix may not be diagonalizable, may have complex eigenvalues, must
be square. The SVD has \emph{none} of these defects. The trick is to launder $A$ through the
one matrix that is always perfectly behaved: $A\T A$ is symmetric and positive semidefinite, so
the \keyword{spectral theorem} hands us an orthonormal eigenbasis on a silver platter. The
right singular vectors are those eigenvectors; the singular values are square roots of the
eigenvalues.
\end{strand}

\block{Existence, in one construction}

\begin{theorem}[every matrix has an SVD]\label{thm:exist}
Every $A\in\R^{m\times n}$ admits a singular value decomposition $A=U\Sig V\T$.
\end{theorem}
\begin{proof}
Skeleton (carry it out as you read):
\begin{enumerate}[leftmargin=2em,itemsep=1pt]
  \item $A\T A$ is symmetric PSD, so by the spectral theorem $A\T A=V\Lambda V\T$ with $V$
        orthonormal and $\Lambda=\diag(\lambda_1\ge\cdots\ge\lambda_n\ge 0)$.
  \item Set $\sigma_i=\sqrt{\lambda_i}$. For each $i$ with $\sigma_i>0$, define
        $u_i=\dfrac{1}{\sigma_i}\,A v_i$. \TODO{check $\{u_i\}$ is orthonormal, using $v_i\T A\T A v_j=\lambda_j\,v_i\T v_j$}
  \item Extend $\{u_i\}$ to an orthonormal basis $U$ of $\R^m$. \TODO{verify $A=\sum_i\sigma_i u_i v_i\T = U\Sig V\T$ by checking both sides agree on the basis $\{v_j\}$}
\end{enumerate}
\end{proof}

\recall{Mirror route: $AA\T=U\Sig\Sig\T U\T$ gives the \emph{left} vectors $u_i$ directly. $A\T A$ and $AA\T$ share the nonzero eigenvalues $\sigma_i^2$.}

\begin{remark}[how unique?]
The singular \emph{values} $\sigma_i$ are unique. The singular \emph{vectors} are not: each
$(u_i,v_i)$ may be flipped in sign together, and within a repeated singular value you may rotate
the corresponding vectors. So $\Sig$ is rigid; $U,V$ have a little freedom.
\end{remark}

\begin{yourturn}
Compute the SVD of $\displaystyle A=\begin{pmatrix}0&2\\3&0\end{pmatrix}$ by the recipe above.
\[
  A\T A=\fillin[2.5cm],\quad \text{eigenvalues } \fillin[1.6cm]\Rightarrow \sigma_1,\sigma_2=\fillin[1.4cm],
\]
\[
  v_1=\fillin[1cm],\ v_2=\fillin[1cm],\qquad
  u_i=\tfrac{1}{\sigma_i}Av_i:\ u_1=\fillin[1cm],\ u_2=\fillin[1cm].
\]
\recall{$A\T A=\diag(9,4)$, so $\sigma=3,2$, $v_i=e_i$, and $u_1=e_2,u_2=e_1$. Notice $U\ne V$: this map is a reflection-and-stretch, not a symmetric one.}
\end{yourturn}

\begin{strand}
When $A$ \emph{is} symmetric PSD, the SVD and the eigendecomposition coincide ($U=V$,
$\sigma_i=\lambda_i$). For symmetric but indefinite $A$, $\sigma_i=|\lambda_i|$ and a sign moves
into $U$. The SVD is the eigendecomposition's well-behaved cousin.
\end{strand}

% ===========================================================================
\section{The four subspaces, in one breath}\label{sec:subspaces}
\warmth{0}\quad\whisper{$U$ splits the output space, $V$ splits the input space, along the rank.}

\begin{strand}
Cut the singular vectors at the rank $r$ (the number of nonzero $\sigma_i$). The first $r$
columns of $V$ and the last $n-r$ split the input $\R^n$ into two orthogonal pieces; the first
$r$ columns of $U$ and the last $m-r$ split the output $\R^m$. Those four pieces are exactly the
\keyword{four fundamental subspaces} of $A$, and the SVD names all of them at once.
\end{strand}

\begin{theorem}[the four subspaces]\label{thm:foursub}
Let $A=U\Sig V\T$ have rank $r$ (so $\sigma_1,\dots,\sigma_r>0$ and the rest vanish). Then:
\end{theorem}
\noindent\begin{tabularx}{\linewidth}{@{}l l L@{}}
\toprule
{\color{indigo}subspace} & {\color{indigo}lives in} & \multicolumn{1}{l}{\color{indigo}is the span of \dots \;(dimension)}\\
\midrule
column space $\col(A)$        & $\R^m$ & $u_1,\dots,u_r$ \;($\dim=\fillin[0.8cm]$)\\
left null space $\Null(A\T)$  & $\R^m$ & $u_{r+1},\dots,u_m$ \;($\dim=m-r$)\\
row space $\col(A\T)$         & $\R^n$ & $\fillin[2.5cm]$ \;($\dim=r$)\\
null space $\Null(A)$         & $\R^n$ & $v_{r+1},\dots,v_n$ \;($\dim=\fillin[1.2cm]$)\\
\bottomrule
\end{tabularx}

\begin{strand}
The picture behind the table: $A$ sends the row space \emph{isomorphically} onto the column
space, stretching the $i$-th axis by $\sigma_i$, and it \emph{annihilates} the null space. Nothing
else happens. Rank--nullity ($\dim\Null(A)=n-r$) and the orthogonal splittings
$\col(A)\perp\Null(A\T)$ in $\R^m$, $\col(A\T)\perp\Null(A)$ in $\R^n$ are now visibly true,
because the columns of an orthogonal matrix are orthonormal.
\end{strand}

\recall{Rank, finally honest: $\rank(A)=$ the number of nonzero singular values. Tiny-but-nonzero $\sigma_i$ are why ``numerical rank'' is a threshold on $\Sig$, not a yes/no.}

\begin{yourturn}
Take $\displaystyle A=\begin{pmatrix}1&1\\1&1\end{pmatrix}$, with SVD given by $\sigma_1=2,\ \sigma_2=0$,
$\,u_1=v_1=\tfrac1{\sqrt2}(1,1)$, $\,v_2=\tfrac1{\sqrt2}(1,-1)$. Read off:
\[
  \rank(A)=\fillin[0.8cm],\quad \col(A)=\operatorname{span}\fillin[2cm],\quad
  \Null(A)=\operatorname{span}\fillin[2cm].
\]
\recall{Rank $1$; column space $=\operatorname{span}\{(1,1)\}$; null space $=\operatorname{span}\{(1,-1)\}$. The map flattens the plane onto a line.}
\end{yourturn}

\begin{strand}
Pseudoinverse preview: because $A$ is a clean isomorphism row space $\to$ column space, it has an
obvious inverse \emph{there} (divide by $\sigma_i$) and zero elsewhere. That is the Moore--Penrose
pseudoinverse of §\ref{sec:faces}.
\end{strand}

% ===========================================================================
\section{The best simple picture}\label{sec:eckart}
\warmth{0}\quad\whisper{Keep the top few singular triples and you have the best low-rank copy of $A$. Provably.}

\begin{strand}
Here is why the SVD is everywhere in data science. Write $A$ as a sum of rank-one layers,
\[
  A=\sum_{i=1}^{r}\sigma_i\,u_i v_i\T ,
\]
ordered by importance $\sigma_1\ge\sigma_2\ge\cdots$. Keep only the top $k$ layers and you get
$A_k=\sum_{i\le k}\sigma_i u_i v_i\T$, the \keyword{truncated SVD}. The miracle: among \emph{all}
matrices of rank $\le k$, $A_k$ is the closest one to $A$. Compression, denoising, PCA,
recommender systems, are all this one theorem.
\end{strand}

\begin{theorem}[Eckart--Young--Mirsky]\label{thm:ey}
For every $B$ with $\rank(B)\le k$,
\[
  \norm{A-B}\;\ge\;\norm{A-A_k},
\]
in both the operator and the Frobenius norm. The optimal error is read off the discarded
singular values:
\[
  \norm{A-A_k}_2=\sigma_{k+1},\qquad
  \norm{A-A_k}_F=\sqrt{\sigma_{k+1}^2+\cdots+\sigma_r^2}.
\]
\end{theorem}
\begin{proof}
Two halves:
\begin{enumerate}[leftmargin=2em,itemsep=1pt]
  \item \emph{The error of $A_k$.} $A-A_k=\sum_{i>k}\sigma_i u_i v_i\T$ is itself in SVD form, so
        its singular values are $\sigma_{k+1},\sigma_{k+2},\dots$. \TODO{read $\norm{\cdot}_2$ and $\norm{\cdot}_F$ off those}
  \item \emph{Nobody does better.} Let $\rank(B)\le k$, so $\Null(B)$ has dimension $\ge n-k$.
        It meets $\operatorname{span}\{v_1,\dots,v_{k+1}\}$ (dimension $k+1$) in some unit vector $w$.
        \TODO{show $Bw=0$ and $\norm{Aw}\ge\sigma_{k+1}$, hence $\norm{A-B}_2\ge\norm{(A-B)w}=\norm{Aw}\ge\sigma_{k+1}$}
\end{enumerate}
\end{proof}

\begin{strand}
The reason this works is the ordering. Total ``energy'' $\norm{A}_F^2=\sum_i\sigma_i^2$ piles up in
the largest few singular values, so throwing away the small ones costs almost nothing. A picture,
a covariance matrix, a ratings table: each is usually \emph{approximately} low rank, meaning its
$\sigma_i$ decay fast, and the truncated SVD is a faithful thumbnail at a fraction of the storage.
\end{strand}

\begin{yourturn}
Suppose $A$ ($4\times4$) has singular values $\sigma=(10,\,6,\,2,\,1)$. For the rank-$2$
truncation $A_2$:
\[
  \norm{A-A_2}_2=\fillin[1cm],\qquad
  \norm{A-A_2}_F=\fillin[2cm],\qquad
  \frac{\norm{A_2}_F^2}{\norm{A}_F^2}=\fillin[2.2cm].
\]
\recall{Spectral error $=\sigma_3=2$; Frobenius error $=\sqrt{2^2+1^2}=\sqrt5$; energy kept $=\frac{10^2+6^2}{10^2+6^2+2^2+1^2}=\frac{136}{141}\approx 96.5\%$. Two numbers carry the matrix.}
\end{yourturn}

\begin{strand}
Caveat worth knowing: Eckart--Young is for \emph{unitarily invariant} norms (operator, Frobenius,
all Schatten-$p$). For other measures (entrywise $\ell_1$, nonnegative or sparse factors), best
low-rank approximation can be NP-hard. The SVD's optimality is special, not universal.
\end{strand}

% ===========================================================================
\section{Faces of the SVD}\label{sec:faces}
\warmth{0}\quad\whisper{One factorization, a dozen classical tools. Each is a small edit to $\Sig$.}

\begin{strand}
With $A=U\Sig V\T$ in hand, a surprising number of textbook objects are just $\Sig$, lightly
edited. Invert the nonzero singular values and you get the pseudoinverse; take their ratio and
you get the condition number; throw $\Sig$ away entirely and keep $UV\T$ and you get the nearest
orthogonal matrix. Here is the keyring.
\end{strand}

\block{The keyring: each tool is $\Sig$, edited}
\noindent\begin{tabularx}{\linewidth}{@{}l l L@{}}
\toprule
{\color{indigo}tool} & {\color{indigo}formula} & \multicolumn{1}{l}{\color{indigo}what it does}\\
\midrule
pseudoinverse $A^{+}$      & $V\Sig^{+}U\T$              & least squares: $x^\star=A^{+}b$ minimizes $\norm{Ax-b}$\\
condition number $\kappa_2(A)$ & $\fillin[1.8cm]$        & how much solving $Ax=b$ amplifies error\\
polar factor $Q$           & $UV\T$                     & the orthogonal matrix \emph{nearest} to $A$\\
PCA directions             & top columns of $V$         & principal axes; $\sigma_i^2\propto$ variance captured\\
operator / Frobenius norm  & $\sigma_1\ \ /\ \sqrt{\textstyle\sum\sigma_i^2}$ & size of $A$ as a map / as a vector\\
best rank-$k$ copy         & truncate $\Sig$ to $k$     & §\ref{sec:eckart}\\
\bottomrule
\end{tabularx}

\begin{definition}[Moore--Penrose pseudoinverse]\label{def:pinv}
$A^{+}=V\Sig^{+}U\T$, where $\Sig^{+}$ replaces each nonzero $\sigma_i$ by $\fillin[1.2cm]$ and
keeps the zeros (and transposes the shape). It is the honest inverse on the part of the space
$A$ handles cleanly: it inverts the stretch on the row space and sends the rest to $0$.
\end{definition}

\recall{When $A$ is square and invertible, $\Sig^{+}=\Sig^{-1}$ and $A^{+}=A^{-1}$. The pseudoinverse only earns its keep when $A$ is rectangular or rank-deficient.}

\begin{yourturn}
Least squares in one line. To minimize $\norm{Ax-b}_2$ over $x$, project $b$ onto the column
space and undo the stretch:
\[
  x^\star \;=\; A^{+}b \;=\; \sum_{\sigma_i>0}\ \frac{u_i\T b}{\sigma_i}\,v_i .
\]
Where does a small $\sigma_i$ make this dangerous, and which face in the keyring measures exactly
that danger? \fillin[4cm]
\recall{A tiny $\sigma_i$ blows up $u_i\T b/\sigma_i$ — the condition number $\kappa_2=\sigma_1/\sigma_r$ is the size of that hazard. Regularization (ridge / truncated SVD) tames it.}
\end{yourturn}

\block{The one idea, recovered}
\begin{strand}
Look back at the master table on the first page and check that you can now fill every row.
A single factorization, $A=U\Sig V\T$, answered all of it: the rank is the count of nonzero
$\sigma_i$, the norms and condition number are simple functions of them, the four subspaces are
the columns of $U$ and $V$ cut at the rank, the best low-rank copy keeps the top few, and the
pseudoinverse inverts what can be inverted. \keyword{Ask a matrix a question; read the answer off
its singular values.}
\end{strand}

\loose{Where to go next: the SVD is to general matrices what the spectral theorem is to symmetric ones. For \emph{huge} $A$ you never form $A\T A$; randomized SVD and Krylov methods get the top $\sigma_i$ from a few matrix-vector products. And the same idea, lifted to operators, is the start of functional analysis.}


\end{document}
